\documentclass{in1150-innlevering}
\title{Innleveringsoppgave 3 i IN1150}
\author{Eilif Akerjordet}
\date{5. februar 2021}
\begin{document}
\maketitle

\section*{Oppgave 1}
\begin{deloppgaver}
    \item
        Påstand: Alle oddetall større enn 51 er primtall.
        \begin{itemize}
            \item 55 er et oddetall større enn 51.
            \item 5 er hverken 55 eller 1.
            \item $\frac{55}{5} = 11$, som er et heltall.
            \item Påstanden er dermed usann.
        \end{itemize}
    \item
        Påstand: Alle formler for $F$ slik at $\neg F$ er en kontradiksjon er gyldig.
        \begin{itemize}
            \item Vi antar at $F$ er sann.
            \item Dette gjør $\neg F$ usann.
            \item Ved samme mynt antar vi at F er usann.
            \item Dette gjør $\neg F$ sann.
            \item Påstanden er dermed sann.
        \end{itemize}
    \item
        Påstand: Alle tautologier har minst to konnektiver.
        \begin{itemize}
            \item $P \imp P$ er en tautologi.
            \item $P \imp P$ har kun 1 konnektiv.
            \item Påstanden er dermed usann.
        \end{itemize}
    \item
        Påstand: Alle formler for $F$ slik at både $F$ og $\neg F$ er oppfyllbare er kontradiksjoner.
        \begin{itemize}
            \item Anta at $F$ og $(F \imp \neg F)$ er sanne samtidig.
            \item Da finnes det en valuasjon som gjør $F$ sann og $(F \imp \neg F)$ sann.
            \item $\imp$-formler betyr at valuasjonen medfører at $\neg F$ er sann.
            \item Definisjonen av $\neg$ medfører at $F$ er usann.
            \item Dette er ikke mulig, da vi har antatt at valuasjonen gjør $F$ sann.
            \item Påstanden er dermed sann.
        \end{itemize}
\end{deloppgaver}
\section*{Oppgave 2}
    \begin{itemize}
        \item
            Som følger av påstanden valuasjonen $v$ gjør formelen $F$ sann $(b)$ følger påstanden om at det finnes en valuasjon som gjør formelen $F$ sann $(a)$ som en logisk konsekvens.  \item
            Hvis alle valuasjoner gjør formelen $F$ sann $(c)$, følger det som en logisk konsekvens av dette at det finnes ingen valuasjon som gjør formelen $F$ usann $(d)$. Det samme gjelder den andre veien også. Man kan derfor si at disse to påstandene er logiske konsekvenser av hverandre.
    \end{itemize}

\section*{Oppgave 3}
\section*{Oppgave 4}
    \begin{itemize}
        \item
            Anta for motsigelse at påstanden om at $P \imp (Q \imp P)$ er en tautologi ikke holder.
        \item
            Da må det finnes en valuasjon som gjør påstanden usann.
        \item
            Den eneste måten å gjøre denne valuasjonen usann på er ved det tilfellet at $(Q \imp P)$ er usann, og $P$ er usann.
        \item
            For at dette skal holde, må $P$ være sann.
        \item
            Dette blir en kontradiksjon. $P$ kan ikke være sann og usann samtidig.
        \item
            Dermed holder påstanden. $P \imp (Q \imp P)$ er en tautologi.
    \end{itemize}

\section*{Oppgave 5}
    Mengden $\{1,2,3,a,b\}$ har følgende relasjoner:
    \begin{deloppgaver}
        \item Anti-symmetrisk
        \item Refleksiv
        \item Symmetrisk
        \item Transitiv, irrefleksiv
    \end{deloppgaver}

\section*{Oppgave 6}
    \begin{itemize}
        \item $R = \{<17,1>,<15,3>,<15,1>,<12,6>,<12,3>,<12,1>,<7,1>,<6,3>,<6,1>,<3,1>\}$
        \item
            R er anti-symmetrisk, irrefleksiv og transitiv.
    \end{itemize}
\section*{Oppgave 7}
    \begin{deloppgaver}
        \item $R = \{<1,1>,<2,2>,<3,3>,<1,2>,<2,1>,<2,3>,<3,2>\}$
        \item $R = \{<1,2>,<2,1>,<2,3>,<3,2>,<1,3>\}$
        \item $R = \{<1,2>,<2,1>,<2,3>,<3,2>,<1,3>\}$
        \item $R = \{<1,2>,<2,1>,<2,3>,<3,2>\}$
    \end{deloppgaver}

\section*{Oppgave 8}
\end{document}
